% --- 1. INFORMAZIONI GENERALI ---
\section{Informazioni Generali}
\begin{itemize}
    \item \textbf{Luogo/Piattaforma:} Server Discord\textsubscript{G}
    \item \textbf{Data:} 2026-06-16
    \item \textbf{Orario di inizio:} 14:30
    \item \textbf{Orario di fine:} 16:30
    \item \textbf{Responsabile\textsubscript{G}:} Sofia De Blasi
    \item \textbf{Scriba:} Sofia De Blasi
\end{itemize}

\vspace{0.5cm}
\textbf{Partecipanti presenti:}
\begin{table}[ht]
    \centering
    \renewcommand{\arraystretch}{1.2}
    \begin{tabularx}{\textwidth}{X l}
        \toprule
        \textbf{Nome e Cognome} & \textbf{Matricola} \\
        \midrule
        Sofia De Blasi & 2111014 \\
        Roberto Mariano Doroftei & 2111031 \\
        Filippo Idiometri & 2035617 \\
        Francesco Marcon & 2110990 \\
        Armando Moda Scarati & 2082864 \\
        Anton Ricardo Rupi & 2054304 \\
        \bottomrule
    \end{tabularx}
\end{table}

\vspace{0.5cm}
\textbf{Partecipanti assenti:}
\begin{table}[ht]
    \centering
    \renewcommand{\arraystretch}{1.2}
    \begin{tabularx}{\textwidth}{X l}
        \toprule
        \textbf{Nome e Cognome} & \textbf{Matricola} \\
        \midrule
        Nessuno & - \\
        \bottomrule
    \end{tabularx}
\end{table}

% --- 2. ORDINE DEL GIORNO ---
\section{Ordine del Giorno}

Gli argomenti previsti per la riunione odierna sono:

\begin{enumerate}
    \item Sprint\textsubscript{G} retrospective Sprint 10 e conclusione della fase RTB\textsubscript{G}.
    \item Revisione\textsubscript{G} finale della documentazione destinata alla consegna:
    \begin{itemize}
        \item Piano di Progetto\textsubscript{G};
        \item Piano di Qualifica\textsubscript{G};
        \item Norme di Progetto\textsubscript{G};
        \item Glossario;
        \item Analisi dei Requisiti\textsubscript{G};
        \item Stato dell'Arte;
        \item Sito documentale;
        \item Lettera di presentazione\textsubscript{G}.
    \end{itemize}
    \item Verifica del corretto funzionamento del Proof of Concept\textsubscript{G} (PoC\textsubscript{G}).
    \item Revisione delle presentazioni per la revisione RTB.
    \item Definizione delle disponibilità del gruppo per la richiesta di colloquio RTB.
    \item Pianificazione delle attività successive alla richiesta di revisione.
\end{enumerate}


% --- 3. RESOCONTO E DISCUSSIONE ---
\section{Resoconto della Riunione}

\subsection{Sprint Retrospective\textsubscript{G} Sprint 10 e Conclusione della Fase RTB}

Il gruppo ha svolto una retrospettiva conclusiva sulla fase RTB.

È stato evidenziato come le convenzioni recentemente introdotte per la nomenclatura dei commit\textsubscript{G} debbano essere applicate con maggiore attenzione nelle attività future.

Inoltre, il gruppo ha osservato che, durante l'ultimo sprint, la concomitanza con gli impegni della sessione d'esame ha comportato una concentrazione di alcune attività in prossimità delle scadenze. Nonostante ciò, grazie alla collaborazione tra i membri del gruppo e al supporto reciproco nelle attività in ritardo, tutte le consegne previste sono state completate nei tempi stabiliti.

Per le fasi successive è stato concordato di distribuire maggiormente le attività durante l'intero sprint, così da garantire al Verificatore\textsubscript{G} un intervallo temporale adeguato per effettuare le verifiche prima delle scadenze.

\subsection{Revisione Finale della Documentazione RTB}

Il gruppo ha effettuato una revisione conclusiva di tutta la documentazione destinata alla consegna, verificando la correttezza dei collegamenti interni ed esterni, dei riferimenti incrociati, dei registri delle modifiche e della coerenza complessiva dei documenti.

In particolare sono stati revisionati:

\begin{itemize}
\item Piano di Progetto;
\item Piano di Qualifica;
\item Norme di Progetto;
\item Glossario;
\item Analisi dei Requisiti;
\item Stato dell'Arte;
\item Sito documentale;
\item Lettera di presentazione.
\end{itemize}

A seguito della revisione non sono emerse criticità bloccanti e il gruppo ha ritenuto la documentazione pronta per la richiesta di revisione RTB.

\subsection{Verifica del Proof of Concept}

Il gruppo ha verificato il corretto funzionamento del PoC sui diversi ambienti di sviluppo dei membri.

È stato ricordato che l'esecuzione del sistema richiede la presenza di un file \texttt{.env} contenente le credenziali necessarie per l'accesso ai servizi LLM\textsubscript{G} utilizzati dal progetto.

La verifica ha confermato la corretta eseguibilità del PoC da parte dei membri del gruppo.

\subsection{Preparazione delle Presentazioni}

Sono state revisionate sia la presentazione relativa allo Stato dell'Arte e alle tecnologie adottate sia la presentazione generale per la revisione RTB.

Il gruppo ha verificato la correttezza dei contenuti, la coerenza con la documentazione prodotta e la distribuzione degli argomenti tra i vari membri per l'esposizione.

Le presentazioni sono state ritenute adeguate per la revisione.

\subsection{Richiesta della Revisione RTB}

Il gruppo ha predisposto la lettera di presentazione destinata ai docenti responsabili della revisione RTB e ne ha verificato la correttezza, inclusa la validità dei collegamenti ai documenti e al sito documentale.

Al fine di agevolare l'organizzazione del colloquio, i membri del gruppo si sono confrontati sulle rispettive disponibilità, tenendo conto degli esami della sessione estiva, individuando un insieme di date e fasce orarie compatibili con gli impegni di tutti i partecipanti.

Al termine della riunione il gruppo ha approvato l'invio della richiesta di revisione RTB tramite posta elettronica al prof. Cardin.

\subsection{Sprint Planning\textsubscript{G} Sprint 11}

A seguito della conclusione della fase RTB e dell'invio della richiesta di revisione, il gruppo ha pianificato le attività previste per lo Sprint 11.

In particolare sono state individuate le seguenti attività:

\begin{itemize}
\item aggiornamento del Piano di Progetto con il consuntivo\textsubscript{G} dello Sprint 10;
\item aggiornamento del Piano di Qualifica con le misurazioni dello Sprint 10;
\item stesura del verbale interno della riunione odierna;
\item stesura del quarto Diario di Bordo\textsubscript{G};
\item preparazione del gruppo alla revisione RTB.
\end{itemize}

Poiché alla data della riunione non erano ancora noti la data della revisione RTB e il relativo esito, il gruppo ha inoltre definito preventivamente la distribuzione dei ruoli da impiegare nel caso in cui fossero richieste modifiche o integrazioni alla documentazione.

Sono stati individuati:

\begin{itemize}
\item Francesco Marcon e Filippo Idiometri come Analisti;
\item Anton Ricardo Rupi e Roberto Mariano Doroftei come Amministratori;
\item Armando Moda Scarati come Verificatore;
\item Sofia De Blasi come Responsabile.
\end{itemize}

Tale organizzazione consentirà di reagire tempestivamente alle eventuali osservazioni emerse durante la revisione RTB e di assegnare rapidamente le attività correttive necessarie.

% --- 4. DECISIONI PRESE ---
\section{Decisioni Prese}

\renewcommand{\arraystretch}{1.3}
\vspace{-0.3cm}
\noindent
\begin{longtable}{c p{12cm}}
    \toprule
    \textbf{Codice} & \textbf{Decisione} \\
    \midrule
    \endfirsthead

    \toprule
    \textbf{Codice} & \textbf{Decisione} \\
    \midrule
    \endhead

    \midrule
    \multicolumn{2}{r}{\textit{Continua nella pagina successiva}} \\
    \midrule
    \endfoot

    \bottomrule
    \endlastfoot

    \textbf{VI-18.1} & La documentazione prodotta per la revisione RTB è ritenuta completa e pronta per la consegna. \\[4pt] 
    \textbf{VI-18.2} & Le presentazioni per la revisione RTB sono approvate per l'esposizione. \\[4pt] 
    \textbf{VI-18.3} & È approvato l'invio della richiesta di revisione RTB al prof. Cardin. \\[4pt] 
    \textbf{VI-18.4} & Viene avviato lo Sprint 11, comprendente l'aggiornamento dei documenti gestionali (PdP e PdQ) con i dati dello Sprint 10 e la stesura del verbale interno e del quarto Diario di Bordo.\\[4pt]
    \textbf{VI-18.5} & Viene approvata la distribuzione preventiva dei ruoli per la gestione di eventuali attività correttive conseguenti alla revisione RTB. \\
\end{longtable}

% --- 5. AZIONI (TODO) ---
\section{Azioni da Svolgere (To-Do)}

\renewcommand{\arraystretch}{1.3}
\begin{longtable}{p{2.3cm} c p{4.5cm} p{3cm} l}
    \toprule
    \textbf{Codice} & \textbf{Rif. Decisione} & \textbf{Descrizione Task\textsubscript{G}} & \textbf{Assegnatario} & \textbf{Scadenza} \\
    \midrule
    \endfirsthead

    \toprule
    \textbf{Codice} & \textbf{Rif. Decisione} & \textbf{Descrizione Task} & \textbf{Assegnatario} & \textbf{Scadenza} \\
    \midrule
    \endhead

    \midrule
    \multicolumn{5}{r}{\textit{Continua nella pagina successiva}} \\
    \midrule
    \endfoot

    \bottomrule
    \endlastfoot

    \ghissue{235} & VI-18.4 & Aggiornamento del Piano di Progetto con il consuntivo dello Sprint 10 & Anton Ricardo Rupi & 2026-06-23 \\[4pt] 
    \ghissue{236} & VI-18.4 & Aggiornamento del Piano di Qualifica con le misurazioni dello Sprint 10 & Anton Ricardo Rupi & 2026-06-23 \\[4pt]
    \ghissue{245} & VI-18.4 & Aggiornamento del Diario di Bordo con la chiusura della fase RTB & Sofia De Blasi & 2026-06-18 \\[4pt] 
    \ghissue{246} & VI-18.4 & Stesura del verbale interno della riunione odierna & Sofia De Blasi & 2026-06-23 \\[4pt] 
    \ghissue{247} & VI-18.3 & Stesura e invio della richiesta di revisione RTB tramite posta elettronica al prof. Cardin & Sofia De Blasi & 2026-06-16 \\
\end{longtable}

% --- 6. PROSSIMA RIUNIONE ---
\section{Prossima Riunione}

\begin{itemize}
\item \textbf{Data:} da concordare (indicativamente 2026-06-23)
\item \textbf{Ora:} da concordare
\item \textbf{OdG preliminare\textsubscript{G}:}
\begin{itemize}
\item aggiornamento sullo stato della richiesta di revisione RTB;
\item analisi dell'eventuale esito della revisione RTB;
\item pianificazione delle eventuali attività correttive richieste;
\item aggiornamento dei documenti gestionali di fine Sprint 10;
\item review dello Sprint 11;
\item organizzazione delle attività successive alla conclusione della fase RTB.
\end{itemize}
\end{itemize}

